\section{Commitments}\label{sec:prelims-commitments}

Commitment schemes enable a party to commit to a message while keeping it hidden
until it is later revealed.

\paragraph{Syntax.} A commitment scheme consists of algorithms defined as
follows:
%
\begin{description}

  \item[$\ck \sample \ComGen(\secparam, n)$ :] The key generation algorithm
takes the security parameter and a vector length $n$ as input, and outputs a
commitment key $\ck$ supporting messages $\msg = (m_1, \ldots, m_n)$.

  \item[$\com \leftarrow \ComCommit(\ck, \msg, \openinfo)$ :] The commitment
algorithm outputs commitment $\com$ to message $\msg$ using key $\ck$ and
opening information $\openinfo$.

  \item[$\msg \leftarrow \ComOpen(\ck, \com, \openinfo)$ :] The opening
algorithm outputs the message $\msg$ corresponding to commitment $\com$ using
opening information $\openinfo$ and key $\ck$.

\end{description}
%
Commitment schemes satisfy correctness, hiding, and binding, which we define
below. We use the standard syntax and security notions for commitment
schemes~\cite{KatLin14}.

The message is a vector. The tracing key of~\Cref{chap:txid} is a pair of field
elements committed together, and both components must be bound by the same
commitment for the proofs of that chapter to relate them. We write $n = 1$ and
drop the vector notation where a scheme is used on a single message. Where a
commitment is opened inside a zero-knowledge proof, the message and the opening
information are both witnesses and the statement proved is that $\ComOpen$
returns the committed value, which the verifier never sees.

\paragraph{Correctness.} A commitment scheme is correct if, for every key $\ck$
generated by $\ComGen(\secparam, n)$, every message $\msg$, and every
$\openinfo$, the following holds:
%
\begin{equation*}
%
  \ComOpen(\ck, \ComCommit(\ck, \msg, \openinfo), \openinfo) = \msg.
  %
\end{equation*}

\paragraph{Hiding.} Hiding requires that a commitment reveal no information
about its message before it is opened.

\begin{experimentbox}{Hiding}{hiding}

  \begin{enumerate}

    \item The challenger generates $\ck \sample \ComGen(\secparam, n)$ and sends
$\ck$ to $\adv$.

    \item The adversary outputs two messages $(\msg_0, \msg_1)$.

    \item The challenger samples $b \sample \{0,1\}$, computes the challenge
commitment $\com^\star \leftarrow \ComCommit(\ck, \msg_b, \openinfo)$, and sends
$\com^\star$ to $\adv$.

    \item The adversary outputs a bit $b'$.

  \end{enumerate}

  $\adv$ wins the experiment if $b' = b$.

\end{experimentbox}

The hiding advantage of $\adv$ is defined as follows:
%
\begin{equation*}
%
  \Adv_{\Com,\adv}^{\mathsf{hiding}}(\secparam) = \left| \Pr[b'=b]-\frac{1}{2}
\right|.
%
\end{equation*}
%
We say the commitment scheme $\Com$ is hiding if this advantage is small for
every admissible adversary $\adv$.

\paragraph{Binding.} We also require commitment schemes to satisfy the binding
property, which prevents an adversary from opening the same commitment to two
different messages.

\begin{experimentbox}{Binding}{binding}

  \begin{enumerate}

    \item The challenger generates $\ck \sample \ComGen(\secparam, n)$ and sends
$\ck$ to $\adv$.

    \item The adversary outputs $(\com, \msg, \msg', \openinfo, \openinfo')$
where $\msg \neq \msg'$.

  \end{enumerate}

  $\adv$ wins the experiment if the following hold:
  %
  \begin{equation*}
  %
    \ComOpen(\ck, \com, \openinfo) = \msg \ \wedge\ \ComOpen(\ck, \com,
\openinfo') = \msg'.
  %
  \end{equation*}

\end{experimentbox}

The binding advantage of $\adv$ is defined as follows:
%
\begin{equation*}
%
  \Adv^{\mathsf{binding}}_{\Com,\adv}(\secparam) = \Pr[\adv\text{ wins}].
%
\end{equation*}
%
A commitment scheme is binding if this advantage is small for every admissible
adversary $\adv$.

\paragraph{Who generates the key.} Both experiments have the challenger generate
$\ck$, which is the standard convention and is harmless as long as the key
reaches the adversary as a value it did not choose. It is not harmless in
deployment. Binding of the algebraic commitments we use rests on no party
knowing the discrete logarithm relations among the group elements of $\ck$, so a
party that samples $\ck$ itself is not bound by commitments under it, and hiding
correspondingly fails against a party that samples the randomness. This matters
wherever the same group elements serve two roles at once. In the instantiation
of~\Cref{sec:txid-instantiation} the commitment key and the registration
authority's signing key share their generators, and binding is required to hold
\emph{against} that authority, so the generators are derived by hash-to-curve
from a public seed fixed in $\crs$ and are an input to key generation rather
than an output of it. \Cref{sec:offline-instantiation} does the same. Where a
scheme is used this way we say so at the point of use, and the security
statements there are made against a key no participant chose.
